\documentclass[10pt,a4paper]{amsart}
\usepackage[margin=19mm]{geometry}
\usepackage{amsmath,amssymb,mathtools}
\usepackage[scr=rsfs,cal=euler]{mathalfa}
\usepackage{xfrac}
\usepackage[unicode,hidelinks,bookmarks=false]{hyperref}
\newcommand{\ie}{i.e.\ }
\newcommand{\cf}{cf.\ }
\newcommand{\resp}{respectively\ }
\newcommand{\Subsection}{\subsection}
\providecommand{\nobreakdash}{\nobreak\hbox{-}}
\setlength{\emergencystretch}{2em}
\multlinegap0pt
\usepackage{bm}
\usepackage{bbm}
\usepackage{dsfont}
\usepackage{xspace}
\usepackage{cancel}
\usepackage{mathtools}
\usepackage{amsthm}
\makeatletter
\@ifundefined{proposition}{\newtheorem{proposition}{Proposition}}{}
\makeatother
\title[Envelope of Holomorphy of Matsuki Orbits in Complex Grassman]{Envelope of Holomorphy of Matsuki Orbits in Complex Grassmannians}
\author{Irfan Ullah}
\date{Source: arXiv:2509.17001v2 (CC BY 4.0)}
\begin{document}
\maketitle

\noindent\emph{Source and licence.} Envelope of Holomorphy of Matsuki Orbits in Complex Grassmannians. Version arXiv:2509.17001v2 (CC BY 4.0), distributed under the Creative Commons Attribution 4.0 International licence, as recorded in the arXiv licence metadata for this exact version and shown on the abstract page https://arxiv.org/abs/2509.17001v2. Full rights notice: licenses/arxiv-2509.17001-CC-BY-4.0.txt.\par\smallskip

\noindent\textit{Editorial scope.} Counted unit: the Proposition whose source label is \texttt{None} in the article (source0.tex lines 229--243, source label \texttt{None}), delivered together with the article's own complete original proof of it (source0.tex lines 245--260, one \texttt{proof} environment of 16 lines). No mathematics is written, repaired or re-derived: statement and proof are the article's own text line for line. Required context retained uncounted: none. Every definition, display and label of those lines is retained unchanged. Omitted from this package and not redistributed: 1--228, 244--244, 261--652. Nothing inside a retained excerpt is omitted or altered: every retained body is delivered exactly as the article's lines are written, so no omission receipt of any kind accompanies this package. Citations: the retained excerpts carry no citation command at all, so no bibliography entry of the article is transcribed and no reference is redistributed. Reference adaptation: none was needed and none was made; every reference inside the delivered unit resolves inside it, so no unresolved reference or citation is produced. Printed slips kept literal: no printed word, symbol, hypothesis or number of the counted statement or of its proof was altered. Layout and class: the article's own class is \texttt{article} with packages including amsmath, amsthm, amssymb, mathrsfs, authblk, lipsum, biblatex; the wrapper is the lane's pinned \texttt{amsart} editorial wrapper at 10pt on A4, which restates the theorem-like environments the retained text needs and adds the article's own macro definitions that the retained text needs (none), with extra wrapper packages (bm, bbm, dsfont, xspace, cancel, mathtools). The delivered numbering is the unit's own. Assets: no retained span contains a figure environment, a graphics-inclusion command, a PDF-page inclusion, a table or a TikZ picture; no image file is copied into this package, the package declares no asset dependency and no image file is redistributed with it. Licence: version arXiv:2509.17001v2 of this article is distributed under the Creative Commons Attribution 4.0 International licence, as recorded in the arXiv licence metadata for this exact version and shown on the abstract page https://arxiv.org/abs/2509.17001v2. A full rights notice is delivered at licenses/arxiv-2509.17001-CC-BY-4.0.txt. Characters: every retained excerpt is pure ASCII, so the delivered engine, XeLaTeX with its default Latin Modern text font, reports no missing character.
\begin{proposition}
For integers $\ell,m \geq 0$ with $\ell+m \leq k$, the $K$-orbit
\[
O_{\ell,m} := \{ W \in \mathrm{Gr}_k(V) : 
\dim(W \cap E_+) = \ell, \ \dim(W \cap E_-) = m \}
\]
is a smooth complex submanifold of $\mathrm{Gr}_k(V)$ and admits the structure 
of a holomorphic fiber bundle
\[
O_{\ell,m} \;\cong\; \mathrm{Gr}_\ell(E_+) \times \mathrm{Gr}_m(E_-) 
\;\ltimes\; O^{(r)}_{0,0},
\]
where $r = k-\ell-m$ and $O^{(r)}_{0,0}$ is the open $K$-orbit of $r$-planes 
in $E_+^\perp \oplus E_-^\perp$. 
\end{proposition}

\begin{proof}[Idea of proof]
Given $W \in O_{\ell,m}$, write 
$W = (W \cap E_+) \oplus (W \cap E_-) \oplus W_{\mathrm{comp}}$ 
with $\dim W_{\mathrm{comp}} = r$. 
By construction $W_{\mathrm{comp}}$ intersects $E_\pm$ trivially, hence 
belongs to $O^{(r)}_{0,0}$. 
The natural projection
\[
\pi : O_{\ell,m} \to \mathrm{Gr}_\ell(E_+) \times \mathrm{Gr}_m(E_-), 
\quad W \mapsto (W \cap E_+, W \cap E_-),
\]
is $K$-equivariant with fiber $O^{(r)}_{0,0}$, yielding the desired 
fiber bundle structure. 
Holomorphicity follows since both base and fiber are complex manifolds, 
and local trivializations are induced from charts on $\mathrm{Gr}_k(V)$. 
\end{proof}

\end{document}
